\documentstyle[%


righttag%


,11pt%


,amssymb%


,russian%               remove in Eng.


,izvru%                 Set izven% in Eng.


]{amsart}





%


\newcommand{\tts}[1]{}





%\hoffset=-15mm%                  For Eng.


\setcounter{page}{23}


\god{2002}


\nomer{5~(480)} %                        Remove in Eng.


%\nomer{Vol.\,46, No.\,5}%               Set in Eng.


%\str{\pageref{begin}--\pageref{end}}%  Set in Eng.


\udk{517.955}





\author{\it В.И.\,Жегалов, А.Н.\,Миронов}


% NB:   ^^ remove in Eng.


\title{О задачах Коши для двух уравнений\\


в частных производных}





\date{}


%\pagestyle{myheadings}%         Set in Eng.


\pagestyle{plain} %              Remove in Eng.


\begin{document}


%\thispagestyle{plain}%          Set in Eng.


\maketitle


\label{begin}











Речь пойдет об уравнении


\begin{equation}


L(u)\equiv u_{xxy}+au_{xx}+bu_{xy}+cu_{x}+du_{y}+eu=f


\end{equation}


и его трехмерном аналоге. Уравнение (1), имеющее приложения, в


частности, в биологии ([1], с.\,261), исследовалось, например, в работах


[2]--[8]. В основном изучалась задача Гурса


и близкие к ней. В данной работе в терминах функций Римана строятся


формулы решения задачи Коши.


\medskip





{\bf 1.} Будем предполагать, что в рассматриваемой области гладкость


коэффициентов определяется включениями $a,\dotsc$, $e\in C^{2}$, $f\in C$.


Следуя работе [8], определим функцию Римана $R(x,y,\xi,\eta)$


как решение интегрального уравнения


\begin{multline}


v(x,y)-\int_{\eta}^{y} a(x,\beta) v(x,\beta) d\beta -


\int_{\xi}^{x} [b(\alpha,y) -(x-\alpha)d(\alpha,y)]v(\alpha,y)d\alpha+\\


+\int_{\xi}^{x} \int_{\eta}^{y} [c(\alpha,\beta) -(x-\alpha)e(\alpha,\beta)]


v(\alpha,\beta) d\beta\, d\alpha=1.


\end{multline}


Решение (2) существует и единственно ([9], cc.\,154, 164).





Справедливо тождество [8]:


\begin{gather}


(uR)_{xxy}\equiv RL(u)+(Mu)_{xy}+(Nu)_{xx}-(Pu)_{x}-(Qu)_{y}+


[u_{y}R_{x}+u(aR)_{x}]_{x}, \\


%


M=R_{x}-bR,\quad N=R_{y}-aR,\quad P=R_{xy}-(aR)_{x}-(bR)_{y}+cR,\notag \\


%


Q=R_{xx}-(bR)_{x}+dR,\notag


\end{gather}


где $R$ зависит от $(x,y,\xi,\eta)$, а коэффициенты $a,\dotsc$, $d$ ---


от $(x,y)$. Из (2) следует, что


\begin{equation}


\begin{array}{c}


M(x,y,x,y)\equiv N(x,y,x,\eta)\equiv P(x,y,x,\eta)\equiv


Q(x,y,\xi,y)\equiv 0,\\[2mm]


R(x,y,x,y)\equiv 1.


\end{array}


\end{equation}


Запишем (3) несколько иначе


\begin{gather}


RL(u)\equiv\frac{\partial S}{\partial x}+\frac{\partial T}{\partial y},\\


%


S=\frac{1}{2}(uR)_{xy}


-\frac{1}{2}[u(2R_{x}-bR)]_{y}+[u(R_{y}-aR)]_{x}+


u[R_{xy}-(aR)_{x}-(bR)_{y}+cR],\notag \\


%


T=\frac{1}{2}(uR)_{xx}-\frac{1}{2}[u(2R_{x}-bR)]_{x}


+u[R_{xx}-(bR)_{x}+dR].\notag


\end{gather}





Пусть $D$ --- треугольная область плоскости $(\xi,\eta)$, ограниченная


характеристиками $\xi=x_{0}$, $\eta=y_{0}$, $x_{0}>0$, $y_{0}>0$, и


отрезком кривой $\Sigma$: $\eta=\sigma(\xi)$, $\sigma'(\xi)<0$,


класса $C^{2}$. Для определенности полагаем $y_{0}=\sigma(0)$,


$\sigma(x_{0})=0$. Сформулируем задачу Коши для уравнения (1): найти


функцию $u\in C^{2}(D\cup\Sigma)\cap C^{(2,1)}(D)$, удовлетворяющую уравнению


(1) и граничным условиям


\begin{equation}


u\big|_{\Sigma}=u_{0}(\xi),\quad


\frac{\partial u}{\partial \bold n}\Big|_{\Sigma}=u_{1}(\xi),


\quad


\frac{\partial^2 u}{\partial \bold n^2}\Big|_{\Sigma}=u_{2}(\xi).


\end{equation}


Здесь $\bold n$ --- единичный вектор внешней нормали,


$\bold n =(\sigma',-1)/\Delta$, $\Delta=\sqrt{1+{\sigma'}^{2}(x)}$


([10], c.\,254), $u_{0}\in C^{2}$, $u_{1}\in C^{1}$, $u_{2}\in C$.


Класс $C^{(k,l)}$ означает существование и непрерывность всех производных


$\partial^{r+s}/\partial x^{r} \partial y^{s}$


($r=0,\dotsc,k$; $s=0,\dotsc,l$).





Рассмотрим точку $(x,y)$ из $D$. Пусть $y_{1}=\sigma(x)$,


$y=\sigma(x_{1})$; $D_{xy}$ и $\Sigma_{xy}$ --- части области $D$ и


кривой $\Sigma$ соответственно, лежащие между характеристиками


$\xi=x$, $\eta=y$. Заменив в (5) переменные $\xi$ на $x$, $\eta$


на $y$, проинтегрируем (5) по $(\xi,\eta)$ по области $D_{xy}$.


Использовав формулу Грина ([11], c.\,236), получим


\begin{equation}


\iint\limits_{D_{xy}} Rfd\xi\, d\eta =


\int_{x_{1}}^{x} T\big|_{\eta=y} d\xi+


\int_{y_{1}}^{y} S\big|_{\xi=x} d\eta


+\int\limits_{\Sigma_{xy}} S d\eta - T d\xi .


\end{equation}


Учитывая тождества (4), после очевидных преобразований приведем (7) к


виду


\begin{multline}


\iint\limits_{D_{xy}} Rfd\xi\, d\eta =


u_{x}(x,y)-\frac{1}{2}(uR)_{\xi}(x_{1},y,x,y)-


\frac{1}{2}(uR)_{\xi}(x,y_{1},x,y)+ \\


+\frac{1}{2}[u(2R_{\xi}-bR)](x_{1},y,x,y)+


\frac{1}{2}[u(2R_{\xi}-bR)](x,y_{1},x,y)


+\int\limits_{\Sigma_{xy}} S d\eta - T d\xi .


\end{multline}


Переписав (8) в виде


\begin{equation}


u_{x}(x,y)=F(x,y)


\end{equation}


и проинтегрировав (9) по $x$, получим решение задачи Коши


\begin{equation}


u(x,y)=u(x_{1},y)+\int_{x_{1}}^{x} F(\alpha,y)d\alpha.


\end{equation}





Формула (10) содержит заданные на $\Sigma$ значения $u$, $u_{x}$,


$u_{y}$, $u_{xx}$, $u_{xy}$. Покажем, что эти значения можно


определить из (6). Действительно, из (6) получаем


\begin{equation}


\begin{split}


\frac{\partial u}{\partial x}+\sigma'\frac{\partial u}{\partial y}&=


u'_{0},\\


\frac{\sigma'}{\Delta}\frac{\partial u}{\partial x}


-\frac{1}{\Delta}\frac{\partial u}{\partial y}&=


u_{1},\\


\sigma''\frac{\partial u}{\partial y}+


\frac{\partial^2 u}{\partial x^2}+


2\sigma'\frac{\partial^2 u}{\partial x \partial y}+


{\sigma'}^{2}\frac{\partial^2 u}{\partial y^2}&=


u''_{0},\\


\frac{\sigma''\Delta-\sigma'\Delta'}{\Delta^2}


\frac{\partial u}{\partial x}


-\frac{\Delta'}{\Delta}\frac{\partial u}{\partial y}


+\frac{\sigma'}{\Delta}\frac{\partial^2 u}{\partial x^2}+


\frac{\sigma'-1}{\Delta}\frac{\partial^2 u}{\partial x \partial y}-


\frac{\sigma'}{\Delta}\frac{\partial^2 u}{\partial y^2}&=


u'_{1},\\


\frac{{\sigma'}^2}{\Delta^2}\frac{\partial^2 u}{\partial x^2}


-\frac{2\sigma'}{\Delta^2}\frac{\partial^2 u}{\partial x \partial y}


+\frac{1}{\Delta^2}\frac{\partial^2 u}{\partial y^2}&=u_{2},


\end{split}


\end{equation}


где все значения берутся на кривой $\Sigma$.


Определитель системы (11)


$$


\frac{1}{\Delta^4}


\begin{vmatrix}


1 & \sigma' & 0 & 0 & 0 \\


\sigma' & -1 & 0 & 0 & 0 \\


0 & \sigma'' & 1 & 2\sigma' & {\sigma'}^2 \\


\dfrac{\sigma''\Delta-\sigma'\Delta'}{\Delta} &


\dfrac{\Delta'}{\Delta} &


\sigma' & {\sigma'}^2 -1 & -\sigma'\\


0 & 0 & {\sigma'}^2 & -2\sigma' & 1


\end{vmatrix}


=\frac{(1+{\sigma'}^2)^4}{\Delta^{4}}=(1+{\sigma'}^2)^2>0,


$$


следовательно, по условиям (6) определяются все требуемые для (10)


функции.


\medskip





{\bf 2.} Рассмотрим теперь трехмерный аналог уравнения (1)


\begin{multline}


L(u)\equiv u_{xxyz}+a_{210}u_{xxy}+a_{201}u_{xxz}+a_{111}u_{xyz}


+a_{200}u_{xx}+a_{110}u_{xy}+\\


+a_{101}u_{xz}+a_{011}u_{yz}+a_{100}u_{x}+a_{010}u_{y}


+a_{001}u_{z}+a_{000}u=f,


\end{multline}


где $a_{klm}\in C^{3}$, $f\in C$.





Развиваем методику из [8]. А именно, назовем функцией Римана


$R(x,y,z,\xi,\eta,\zeta)$ решение интегрального уравнения


\begin{multline}


v(x,y,z)-\int_{\zeta}^{z} a_{210}(x,y,\gamma)v(x,y,\gamma) d\gamma


-\int_{\eta}^{y} a_{201}(x,\beta,z)v(x,\beta,z) d\beta- \\


%


-\int_{\xi}^{x} [a_{111}(\alpha,y,z)-(x-\alpha)a_{011}(\alpha,y,z)]


v(\alpha,y,z) d\alpha+ 


\int_{\eta}^{y} \int_{\zeta}^{z} a_{200}(x,\beta,\gamma)


v(x,\beta,\gamma) d\gamma\, d\beta+\\


%


+\int_{\xi}^{x} \int_{\zeta}^{z} [a_{110}(\alpha,y,\gamma)


-(x-\alpha)a_{010}(\alpha,y,\gamma)]v(\alpha,y,\gamma) d\gamma\, d\alpha+\\


%


+\int_{\xi}^{x} \int_{\eta}^{y} [a_{101}(\alpha,\beta,z)


-(x-\alpha)a_{001}(\alpha,\beta,z)]v(\alpha,\beta,z) d\beta\, d\alpha- \\


%


-\int_{\xi}^{x} \int_{\eta}^{y} \int_{\zeta}^{z}


[a_{100}(\alpha,\beta,\gamma)


-(x-\alpha)a_{000}(\alpha,\beta,\gamma)]v(\alpha,\beta,\gamma)


d\gamma\, d\beta\, d\alpha=1.


\end{multline}


Решение (13) существует и единственно ([9], cc.\,154, 164).





Далее используются обозначения


$$


P_{1}=R_{z}-a_{210}R,\quad P_{2}=R_{y}-a_{201}R,\quad


P_{3}=R_{x}-a_{111}R, 


$$


\begin{alignat*}{2}


Q_{1}&=R_{yz}-(a_{210}R)_{y}-(a_{201}R)_{z}+a_{200}R, &\quad


Q_{2}&=R_{xz}-(a_{210}R)_{x}-(a_{111}R)_{z}+a_{110}R, \\


%


Q_{3}&=R_{xy}-(a_{210}R)_{x}-(a_{111}R)_{y}+a_{101}R, &\quad


%


Q_{4}&=R_{xx}-(a_{111}R)_{x}+a_{011}R, 


\end{alignat*}


\begin{align*}


S_{1}&=R_{xyz}-(a_{210}R)_{xy}-(a_{201}R)_{xz}-(a_{111}R)_{yz}+


(a_{200}R)_{x}+(a_{110}R)_{y}+(a_{101}R)_{z}-a_{110}R, \\


%


S_{2}&=R_{xxz}-(a_{210}R)_{xx}-(a_{111}R)_{xz}+(a_{110}R)_{x}+


(a_{011}R)_{z}-a_{010}R, \\


%


S_{3}&=R_{xxy}-(a_{210}R)_{xx}-(a_{111}R)_{xy}+(a_{101}R)_{x}+


(a_{011}R)_{y}-a_{001}R. 


\end{align*}


Здесь $R=R(x,y,z,\xi,\eta,\zeta)$, а остальные функции зависят от


$(x,y,z)$.





Из интегрального уравнения (13) могут быть получены тождества


\begin{equation} %14


\begin{array}{c}


P_{1}(x,y,\zeta,x,y,z)\equiv P_{2}(x,\eta,z,x,y,z)\equiv


P_{3}(x,y,z,x,y,z)\equiv \\[2mm]


\equiv Q_{1}(x,\eta,\zeta,x,y,z)\equiv Q_{2}(x,y,\zeta,x,y,z)\equiv


Q_{3}(x,\eta,z,x,y,z)\equiv \\[2mm]


\equiv Q_{4}(\xi,y,z,x,y,z)\equiv S_{1}(x,\eta,\zeta,x,y,z)\equiv


S_{2}(\xi,y,\zeta,x,y,z)\equiv \\[2mm]


\equiv S_{3}(\xi,\eta,z,x,y,z)\equiv 0, \quad


R(x,y,z,x,y,z)\equiv 1.


\end{array}


\end{equation}





Непосредственной проверкой можно убедиться в справедливости тождества


\begin{equation} %15


RL(u)\equiv \frac{\partial W_1}{\partial x}


+\frac{\partial W_2}{\partial y}+\frac{\partial W_3}{\partial z},


\end{equation}


\begin{align*}


W_1&=\frac{1}{3}(Ru)_{xyz}-\frac{1}{2}(P_{1}u)_{xy}


-\frac{1}{2}(P_{2}u)_{xz}+(Q_{1}u)_{x}+


%


\frac{1}{2}((Q_{2}+P_{1x})u)_{y}


+\frac{1}{2}((Q_{3}+P_{2x})u)_{z}


-(S_{1}+Q_{1})u, \\


%


W_2&=\frac{1}{3}(Ru)_{xxz}-\frac{1}{2}((P_{3}+R_{x})u)_{xz}


-\frac{1}{2}(P_{1}u)_{xx}+


%


\frac{1}{2}((Q_{2}+P_{1x})u)_{x}


+\frac{1}{2}(Q_{4}u)_{z}-S_{2}u, \\


%


W_3&=\frac{1}{3}(Ru)_{xxy}-\frac{1}{2}((P_{3}+R_{x})u)_{xy}


-\frac{1}{2}(P_{2}u)_{xx}+


%


\frac{1}{2}((Q_{3}+P_{2x})u)_{x}


+\frac{1}{2}(Q_{4}u)_{y}-S_{3}u.


\end{align*}





В дальнейшем используется схема рассуждений статьи [12].


Пусть $\zeta=\zeta(\xi,\eta)$ --- поверхность класса $C^3$ в пространстве


$(\xi,\eta,\zeta)$. Потребуем, чтобы эта поверхность в


каждой своей точке имела касательную плоскость, не параллельную ни


одной из координатных осей. Для определенности можно положить


$\zeta'_{\xi}<0$, $\zeta'_{\eta}<0$.





Проведем через точку $M(x,y,z)$ плоскости $\xi=x$, $\eta=y$, $\zeta=z$.


Пусть указанные плоскости пересекают поверхность $\zeta=\zeta(\xi,\eta)$


по кривым $BC$, $CA$ и $AB$ соответственно.


Плоскости $\xi=x$, $\eta=y$, $\zeta=z$ и поверхность $\zeta=\zeta(\xi,\eta)$


определяют область $D$, граница которой состоит из двумерных


многообразий $AMC$, $BCM$, $AMB$ и $ABC$. Ориентацию $D$ считаем


положительной (ориентацию в пространстве можно связать с направлением


внешней нормали к границе области).





Задача Коши: найти функцию $u\in C^{3}\cap C^{(2,1,1)}$,


удовлетворяющую уравнению (12) и следующим условиям на поверхности


$\zeta=\zeta(\xi,\eta)$:


\begin{equation} %16


u\big|_{ABC}=u_{0}(\xi,\eta),\quad


\frac{\partial u}{\partial\bold n}\Big|_{ABC}=u_{1}(\xi,\eta),


\quad


\frac{\partial^2 u}{\partial^2 \bold n}\Big|_{ABC}


=u_{2}(\xi,\eta),\quad


\frac{\partial^3 u}{\partial^3 \bold n}\Big|_{ABC}


=u_{3}(\xi,\eta).


\end{equation}


Здесь $\bold n$ --- единичный вектор внешней нормали к поверхности


$\zeta=\zeta(\xi,\eta)$, $u_{0}\in C^{3}$, $u_{1}\in C^{2}$,


$u_{2}\in C^{1}$, $u_{3}\in C$.


Класс $C^{(k,l,m)}$ означает существование и непрерывность всех


производных $\partial^{r+s+t}/\partial x^{r} \partial y^{s}


\partial z^{t}$ ($r=0,\dotsc,k$; $s=0,\dotsc,l$; $t=0,\dotsc,m$).





Заменив в (15) переменные $\xi$ на $x$, $\eta$ на $y$, $\zeta$ на


$z$, проинтегрируем (15) по $\xi$, $\eta$, $\zeta$ по области~$D$.


Применяя формулу Гаусса--Остроградского ([11], c.\,241), получим


\begin{equation} %17


\iiint\limits_{D} Rf\,d\xi\, d\eta\, d\zeta=


\iint\limits_{\partial D} W_1\, d\eta \wedge d\zeta


+W_2\, d\zeta \wedge d\xi+W_3\, d\xi \wedge d\eta.


\end{equation}


Здесь знак ``$\wedge$'' обозначает внешнее умножение дифференциальных форм.


Обозначим правую часть (17) через $I$. Заменим интеграл по $\partial D$


суммой интегралов по ее составляющим $ABC$, $AMC$, $BCM$ и $ABM$.


При этом учтем, что в соответствии с (14)


$$


Q_{1}(x,\eta,\zeta,x,y,z)\equiv S_{1}(x,\eta,\zeta,x,y,z)\equiv


S_{2}(\xi,y,\zeta,x,y,z)\equiv S_{3}(\xi,\eta,z,x,y,z)\equiv 0.


$$


Таким образом, получаем


\begin{gather*}


I=\iint\limits_{BCM} \bigg\{


\frac{\partial }{\partial \eta} \bigg[ \frac{1}{6}


(Ru)_{\xi\zeta}-\frac{1}{2} (P_1 u)_\xi +\frac{1}{2} (Q_2+P_{1\xi})u\bigg]+\\


%


+\frac{\partial }{\partial \zeta} \bigg[


\frac{1}{6} (Ru)_{\xi\eta}-\frac{1}{2} (P_2 u)_\xi


+\frac{1}{2} (Q_3+P_{2\xi})u \bigg] \bigg \}\, d\eta\wedge d\zeta + \\


%


+\iint\limits_{AMC} \bigg\{


\frac{\partial }{\partial \xi} \bigg[ \frac{1}{6}


(Ru)_{\xi\zeta}-\frac{1}{4} ((P_3 +R_\xi)u)_\zeta -\frac{1}{2} (P_1 u)_\xi


+\frac{1}{2} (Q_2+P_{1\xi})u \bigg]+ \\


%


+\frac{\partial }{\partial \zeta} \bigg[ \frac{1}{6}


(Ru)_{\xi\xi}-\frac{1}{4} ((P_3 +R_\xi)u)_\xi +\frac{1}{2} Q_4 u


\bigg] \bigg\} \, d\zeta\wedge  d\xi + \\


%


+\iint\limits_{ABM} \bigg \{ \frac{\partial }{\partial \xi} \bigg[


\frac{1}{6} (Ru)_{\xi\eta}-\frac{1}{4}


((P_3 +R_\xi)u)_\eta -\frac{1}{2} (P_2 u)_\xi


+\frac{1}{2} (Q_3 +P_{2\xi})u\bigg] + \\


%


+\frac{\partial }{\partial \eta} \bigg[


\frac{1}{6} (uR)_{\xi\xi} -\frac{1}{4} ((P_3 +R_\xi)u)_\xi +\frac{1}{2} Q_4 u


\bigg] \bigg\} \, d\xi\wedge d\eta + \\


%


+\iint\limits_{ABC} W_1\, d\eta\wedge d\zeta


+W_2\, d\zeta \wedge d\xi + W_3\, d\xi \wedge d\eta.


\end{gather*}





По формуле Грина ([11], c.\,236) интегралы по плоским областям


$AMC$, $BCM$ и $ABM$ сводятся к однократным интегралам по замкнутым


контурам


\begin{gather*}


I=\int\limits_{BCM}


\bigg[ \frac{1}{6} (Ru)_{\xi\zeta}-\frac{1}{2} (P_1 u)_\xi


+\frac{1}{2} (Q_2+P_{1\xi})u \bigg] d\zeta - \\


%


-\bigg[ \frac{1}{6} (Ru)_{\xi\eta}-\frac{1}{2} (P_2 u)_\xi


+\frac{1}{2} (Q_3+P_{2\xi})u \bigg] d\eta + \\


%


+\int\limits_{AMC}


\bigg[ \frac{1}{6} (Ru)_{\xi\xi}-\frac{1}{4}


((P_3 +R_\xi)u)_\xi +\frac{1}{2} Q_4 u \bigg] d\xi - \\


%


-\bigg[ \frac{1}{6} (Ru)_{\xi\zeta}-\frac{1}{4}


((P_3 +R_\xi)u)_\zeta -\frac{1}{2}


(P_1 u)_\xi +\frac{1}{2} (Q_2+P_{1\xi})u \bigg] d\zeta + \\


%


+\int\limits_{ABM} \bigg[


\frac{1}{6} (Ru)_{\xi\eta}-\frac{1}{4}


((P_3 +R_\xi)u)_\eta -\frac{1}{2} (P_2 u)_\xi


+\frac{1}{2} (Q_3 +P_{2\xi})u \bigg] d\eta - \\


%


-\bigg[ \frac{1}{6} (uR)_{\xi\xi} -\frac{1}{4}


((P_3 +R_\xi)u)_\xi +\frac{1}{2} Q_4 u \bigg] d\xi + \\


%


+\iint\limits_{ABC} W_1\, d\eta\wedge d\zeta


+W_2\, d\zeta \wedge d\xi + W_3\, d\xi \wedge d\eta.


\end{gather*}


Заменяем теперь каждый криволинейный интеграл суммой интегралов по


составляющим его контура. Обозначим сумму этих однократных интегралов через


$J$:


\begin{gather*}


J=\int\limits_{CM}


\bigg[ \frac{1}{6} (Ru)_{\xi\zeta}-\frac{1}{2} (P_1 u)_\xi


+\frac{1}{2} (Q_2+P_{1\xi})u \bigg] d\zeta - \\


%


-\int\limits_{MB}


\bigg[ \frac{1}{6} (Ru)_{\xi\eta}-\frac{1}{2} (P_2 u)_\xi


+\frac{1}{2} (Q_3+P_{2\xi})u \bigg] d\eta + \\


%


+\int\limits_{BC}


\bigg[ \frac{1}{6} (Ru)_{\xi\eta}-\frac{1}{2} (P_2 u)_\xi


+\frac{1}{2} (Q_3+P_{2\xi})u \bigg] d\zeta - \\


%


-\bigg[ \frac{1}{6} (Ru)_{\xi\zeta}-\frac{1}{2} (P_1 u)_\xi


+\frac{1}{2} (Q_2+P_{1\xi})u \bigg] d\eta + \\


%


\displaybreak[3]


+\int\limits_{AM}


\bigg[ \frac{1}{6} (Ru)_{\xi\xi}-\frac{1}{4}


((P_3 +R_\xi)u)_\xi +\frac{1}{2} Q_4 u \bigg] d\xi - \\


%


-\int\limits_{MC}


\bigg[ \frac{1}{6} (Ru)_{\xi\zeta}-\frac{1}{4}


((P_3 +R_\xi)u)_\zeta -\frac{1}{2}


(P_1 u)_\xi +\frac{1}{2} (Q_2+P_{1\xi})u \bigg] d\zeta + \\


%


+\int\limits_{CA}


\bigg[ \frac{1}{6} (Ru)_{\xi\xi}-\frac{1}{4}


((P_3 +R_\xi)u)_\xi +\frac{1}{2} Q_4 u\bigg] d\xi - \\


%


-\bigg[ \frac{1}{6} (Ru)_{\xi\zeta}-\frac{1}{4}


((P_3 +R_\xi)u)_\zeta -\frac{1}{2}


(P_1 u)_\xi +\frac{1}{2} (Q_2+P_{1\xi})u \bigg] d\zeta + \\


%


+\int\limits_{BM} \bigg[


\frac{1}{6} (Ru)_{\xi\eta}-\frac{1}{4}


((P_3 +R_\xi)u)_\eta -\frac{1}{2} (P_2 u)_\xi


+\frac{1}{2} (Q_3 +P_{2\xi})u \bigg] d\eta - \\


%


-\int\limits_{MA}


\bigg[ \frac{1}{6} (uR)_{\xi\xi} -\frac{1}{4}


((P_3 +R_\xi)u)_\xi +\frac{1}{2} Q_4 u\bigg] d\xi + \\


%


+\int\limits_{AB} \bigg[


\frac{1}{6} (Ru)_{\xi\eta}-\frac{1}{4}


((P_3 +R_\xi)u)_\eta -\frac{1}{2} (P_2 u)_\xi


+\frac{1}{2} (Q_3 +P_{2\xi})u \bigg] d\eta - \\


%


-\bigg[ \frac{1}{6} (uR)_{\xi\xi} -\frac{1}{4}


((P_3 +R_\xi)u)_\xi +\frac{1}{2} Q_4 u\bigg] d\xi.


\end{gather*}


Учитываем следующие тождества из (14):


\begin{multline*}


P_{1}(x,y,\zeta,x,y,z)\equiv P_{2}(x,\eta,z,x,y,z)\equiv


Q_{1}(x,\eta,\zeta,x,y,z)\equiv\\


\equiv Q_{2}(x,y,\zeta,x,y,z)\equiv Q_{3}(x,\eta,z,x,y,z)


\equiv Q_{4}(\xi,y,z,x,y,z)\equiv 0.


\end{multline*}


Получаем


\begin{gather*}


J=\int\limits_{CM} \frac{1}{6} (Ru)_{\xi\zeta} d\zeta


-\int\limits_{MB} \frac{1}{6} (Ru)_{\xi\eta} d\eta +


\int\limits_{BC}


\bigg[ \frac{1}{6} (Ru)_{\xi\eta}-\frac{1}{2} (P_2 u)_\xi


+\frac{1}{2} (Q_3+P_{2\xi})u \bigg] d\zeta -\\


%


-\bigg[ \frac{1}{6} (Ru)_{\xi\zeta}-\frac{1}{2} (P_1 u)_\xi


+\frac{1}{2} (Q_2+P_{1\xi})u \bigg] d\eta +


\int\limits_{AM}


\bigg[\frac{1}{6}(Ru)_{\xi\xi}-\frac{1}{4} ((P_3 +R_\xi)u)_\xi \bigg] d\xi-\\


%


-\int\limits_{MC}


\bigg[ \frac{1}{6} (Ru)_{\xi\zeta}-\frac{1}{4}


((P_3 +R_\xi)u)_\zeta \bigg] d\zeta +


\int\limits_{CA}


\bigg[ \frac{1}{6} (Ru)_{\xi\xi}-\frac{1}{4}


((P_3 +R_\xi)u)_\xi +\frac{1}{2} Q_4 u


\bigg] d\xi -\\


%


-\bigg[ \frac{1}{6} (Ru)_{\xi\zeta}-\frac{1}{4}


((P_3 +R_\xi)u)_\zeta -\frac{1}{2}


(P_1 u)_\xi +\frac{1}{2} (Q_2+P_{1\xi})u \bigg] d\zeta + \\


%


+\int\limits_{BM} \bigg[


\frac{1}{6} (Ru)_{\xi\eta}-\frac{1}{4} ((P_3 +R_\xi)u)_\eta \bigg] d\eta-


\int\limits_{MA}


\bigg[ \frac{1}{6} (uR)_{\xi\xi} -\frac{1}{4}


((P_3 +R_\xi)u)_\xi \bigg] d\xi + \\


%


+\int\limits_{AB} \bigg[


\frac{1}{6} (Ru)_{\xi\eta}-\frac{1}{4}


((P_3 +R_\xi)u)_\eta -\frac{1}{2} (P_2 u)_\xi


+\frac{1}{2} (Q_3 +P_{2\xi})u \bigg] d\eta - \\


%


-\bigg[ \frac{1}{6} (uR)_{\xi\xi} -\frac{1}{4}


((P_3 +R_\xi)u)_\xi +\frac{1}{2} Q_4 u\bigg] d\xi.


\end{gather*}


Вычисляя в полученной формуле для $J$ интегралы по отрезкам прямых,


запишем (17) в следующем виде:


\begin{multline}


\iiint\limits_{D} Rf\,d\xi\, d\eta\, d\zeta=


u_{\xi}\big|_{M} -\frac{1}{3}\big(


(Ru)_{\xi}\big|_{A}+ (Ru)_{\xi}\big|_{B}+


(Ru)_{\xi}\big|_{C}\big)+ \\


%


+\frac{1}{2}(P_3 +R_\xi)u\big|_{A}


+\frac{1}{4}(P_3 +R_\xi)u\big|_{B}


+\frac{1}{4}(P_3 +R_\xi)u\big|_{C}+ \\


%


+\int\limits_{BC}


\bigg[ \frac{1}{6} (Ru)_{\xi\eta}-\frac{1}{2} (P_2 u)_\xi


+\frac{1}{2} (Q_3+P_{2\xi})u \bigg] d\zeta - \\


%


-\bigg[ \frac{1}{6} (Ru)_{\xi\zeta}-\frac{1}{2} (P_1 u)_\xi


+\frac{1}{2} (Q_2+P_{1\xi})u \bigg] d\eta + \\


%


+\int\limits_{CA}


\bigg[ \frac{1}{6} (Ru)_{\xi\xi}-\frac{1}{4}


((P_3 +R_\xi)u)_\xi +\frac{1}{2} Q_4 u\bigg] d\xi - \\


%


-\bigg[ \frac{1}{6} (Ru)_{\xi\zeta}-\frac{1}{4}


((P_3 +R_\xi)u)_\zeta -\frac{1}{2}


(P_1 u)_\xi +\frac{1}{2} (Q_2+P_{1\xi})u \bigg] d\zeta + \\


%


+\int\limits_{AB} \bigg[


\frac{1}{6} (Ru)_{\xi\eta}-\frac{1}{4}


((P_3 +R_\xi)u)_\eta -\frac{1}{2} (P_2 u)_\xi


+\frac{1}{2} (Q_3 +P_{2\xi})u \bigg] d\eta - \\


%


-\bigg[ \frac{1}{6} (uR)_{\xi\xi} -\frac{1}{4}


((P_3 +R_\xi)u)_\xi +\frac{1}{2} Q_4 u\bigg] d\xi+ \\


%


+\iint\limits_{ABC} W_1\, d\eta\wedge d\zeta


+W_2\, d\zeta \wedge d\xi + W_3\, d\xi \wedge d\eta.


\end{multline}


Запишем (18) в виде


\begin{equation} %19


u_{x}(x,y,z)=F(x,y,z),


\end{equation}


а затем проинтегрируем (19) по $x$. Если точка $A$ имеет координаты


$(x_1,y,z)$, то решение задачи Коши принимает вид


\begin{equation} %20


u(x,y,z)=u(x_{1},y,z)+\int_{x_{1}}^{x} F(\alpha,y,z)d\alpha.


\end{equation}





Все значения функции $u$ и ее производных, входящие в формулу (19),


определяются из условий (16). Это можно показать, рассуждая как в


работе [12]. Пусть $\bold n =(n_1(x,y),n_2(x,y),n_3(x,y))$.


Тогда в криволинейных координатах $(\xi,\eta,\mu)$, связанных с поверхностью


$\zeta=\zeta(\xi,\eta)$, $u=u(\xi+n_1(\xi,\eta)\mu, \eta+n_2(\xi,\eta)\mu,


\zeta(\xi,\eta)+n_3(x,y)\mu)$. Находя на поверхности $\zeta=\zeta(\xi,\eta)$


частные производные функции $u$ по переменным $\xi$, $\eta$, $\mu$,


получаем систему для определения частных производных решения $u$


до третьего порядка включительно на поверхности $\zeta=\zeta(\xi,\eta)$.


Ввиду громоздкости соответствующие выкладки не приводятся.





Заметим, что (10) и (20) играют роль известной формулы Римана


([13], c.\,67, формула (1.169)).





\begin{thebibliography}{99}





\bibitem{1}


Нахушев А.M. {\em Уравнения математической биологии}. -- М.: Высш. школа,


1995. -- 301~с.


\bibitem{2}


Colton D. {\em Pseudoparabolic equations in one space variable}


// J. Different. equations. -- 1972. -- V.\,12. -- \No\,3. --


P.\,559--565.


\bibitem{3}


Rundell W., Stecher M. {\em Remarks 


concerning the support of solutions of pseudoparabolic equation}


// Proc. Amer. Math. Soc. -- 1977. -- V.\,63. -- \No\,1. --


P.\,77--81.


\bibitem{4}


Шхануков М.Х. {\em О некоторых краевых задачах для уравнения третьего порядка,


возникающих при моделировании фильтрации жидкости в пористых средах} //


Дифференц. уравнения. -- 1982. -- Т. 18. -- \No\,4. --


С.\,689--699.


\bibitem{5}


Шхануков М.Х. {\em Об одном методе решения краевых задач для уравнений


третьего порядка} // ДАН CCCP. -- 1982. -- Т.\,265. -- \No\,6. --


С.\,1327--1330.


\bibitem{6}


Солдатов А.П., Шхануков М.Х. {\em Краевые задачи с общим


нелокальным условием А.А.\,Самарского для псевдопараболических уравнений


высокого порядка} // ДАН CCCP. -- 1987. -- Т.\,297. -- \No\,3. --


С.\,547--552.


\bibitem{7}


Джохадзе О.М. {\em Задача типа Дарбу для уравнения третьего порядка с


доминирующими младшими членами} // Дифференц. уравнения. -- 1996. --


Т.\,32. -- \No\,4. -- С.\,523--535.


\bibitem{8}


Жегалов В.И., Уткина Е.А.


{\em Об одном псевдопараболическом уравнении третьего порядка}


// Изв. вузов. Математика. -- 1999. -- \No\,10. -- С.\,73--76.


\bibitem{11}


Мюнтц Г. {\em Интегральные уравнения}. Т.\,1. -- Л.--М.: ГТТИ, 1934. --


330 с.


\bibitem{9}


Владимиров В.С. {\em Уравнения математической физики}. --


2-е изд. -- М.: Наука, 1971. -- 512~с.


\bibitem{10}


Зорич В.А. {\em Математический анализ}. Ч.\,2. -- М.: Наука, 1984. --


640 с.


\bibitem{12}


Севастьянов В.А. {\em Метод Римана для трехмерного


гиперболического уравнения третьего порядка} // Изв. вузов.


Математика. -- 1997. -- \No\,5. -- С.\,69--73.


\bibitem{13}


Бицадзе А.В. {\em Некоторые классы уравнений в частных


производных}. -- М.: Наука, 1981. -- 448~с.





\end{thebibliography}





\vskip 2cm





\post{\it Казанский государственный университет}{\it Поступила}


\post{\it Елабужский государственный}{27.02.2001}


\post{\it педагогический институт}{}





\label{end}


\end{document}





